\RequirePackage[l2tabu,orthodox]{nag}

\documentclass[a0paper,portrait]{baposter}
\usepackage{amsmath}
\usepackage{amssymb}
\usepackage{amsfonts}
\usepackage{mathtools}
\usepackage{microtype}
\usepackage{multicol}
\usepackage{url}

\usepackage{tikz}
\usepackage{pgfplots}
\pgfplotsset{compat=1.13}
\usepackage{cmbright}
\newcommand{\fitCurve}[4]{%
  \pgfplotstablegetelem{#2}{C}\of#1
  \pgfmathsetmacro{#3}{\pgfplotsretval}
  \pgfplotstablegetelem{#2}{a}\of#1
  \pgfmathsetmacro{\a}{\pgfplotsretval}
  \pgfplotstablegetelem{#2}{b}\of#1
  \pgfmathsetmacro{\b}{\pgfplotsretval}
  \addplot[densely dashdotted,domain=\a:\b] {#4};
}
\usepackage{graphicx}
\usepackage{graphbox}
\graphicspath{{./\jobname.figures/}{../../pictures/}}

\usepackage[most]{tcolorbox}
\definecolor{firedrake_red}{RGB}{221,0,0}
\definecolor{firedrake_blue}{RGB}{16,66,122}

\begin{document}

\begin{poster}{
    grid=false,
    columns=6,
    colspacing=0.5em,
    textborder=faded,
    headerColorOne=firedrake_red,
    borderColor=firedrake_red,
    headerFontColor=white,
    headerborder=open,
    headerfont=\bfseries\Large\scshape,
    eyecatcher=true,
    headershape=rounded,
    headershade=plain,
    background=none,
    headerheight=0.09\textheight,
  }
  {
    \makebox[0.3\textwidth][l]{
    \begin{tabular}{r}
      \includegraphics[width=0.15\textwidth]{imperial-two-tone}\\
      \vspace{0.5em}
      \includegraphics[height=0.05\textwidth]{baylor-logo-bw}\\
    \end{tabular}
    }
  } 
  % Title
  {\scshape\Large FInAT is not a tabulator}
  % Authors
  {\small\scshape Lawrence~Mitchell; Mikl\'os~Homolya; David~A.~Ham\\[0.1em]Imperial College London
    \\[0.4em]  
    Robert~C.~Kirby\\[0.1em]Baylor University}
  % University logo
  {
    \makebox[0.3\textwidth][r]{
      \includegraphics[align=c,height=0.05\textheight]{firedrake-word}
    }
  }

  \begin{posterbox}[name=introduction,span=6,column=0,row=0]{FIAT is a
      tabulator}
    \begin{minipage}{0.49\textwidth}
      \subsection*{FIAT}
      \vspace{-0.5em}
      Provides implementations of a wide class of
      finite elements on reference cells. It can tabulate basis
      functions, and their derivatives, at sets of points.
      \vspace{-0.25\baselineskip}
      \begin{align*}
        \Psi_i(\xi_q) &\mapsto \mathbf{\Psi}_{iq}\\
        \hat{\nabla}\Psi_i(\xi_q) &\mapsto D\mathbf{\Psi}_{iq}
      \end{align*}
    \end{minipage}
    \hfill
    \begin{minipage}{0.49\textwidth} 
      \subsection*{FInaT}
      \vspace{-0.5em}
      Provides \emph{symbolic expressions} for the
      evaluation of basis functions at points.  These enable a
      compiler to apply optimisations that depend on \emph{structure}
      in the basis functions.

      We have implemented these optimisations in TSFC, to provide low
      complexity finite element evaluation in Firedrake.
    \end{minipage}
  \end{posterbox}

  \begin{posterbox}[name=structure, column=0, below=introduction, span=3]{Tabulation expressions}
    Arbitrarily structured basis evaluations, composed of six base cases.
    \tcbset{center title, colframe=firedrake_blue,colback=firedrake_blue!10!white,fonttitle=\bfseries}
    \begin{tcbraster}[raster columns=6,raster equal height=rows]
    \begin{tcolorbox}[title=FIAT,left=0mm,right=0mm,top=0mm,bottom=0mm,boxsep=0.5mm,
      raster multicolumn=3, valign upper=center, halign=center]
      Tabulation as a literal array.
      \begin{equation*}
        \Psi_i(\xi_q) \mapsto \mathbf{\Psi}_{iq}
      \end{equation*}
    \end{tcolorbox}
    \begin{tcolorbox}[title=Tensor-valued,left=0mm,right=0mm,top=0mm,bottom=0mm,boxsep=0.5mm,
      raster multicolumn=3, valign upper=center]
      \begin{equation*}
        \Psi_{(\alpha, \nu)\kappa}(\xi_q) \mapsto \Psi_{\alpha}^*(\xi_q) \prod_{i=1}^n \delta_{\nu_i \kappa_i}
      \end{equation*}
    \end{tcolorbox}
    \begin{tcolorbox}[title=Underintegration,left=0mm,right=0mm,top=0mm,bottom=0mm,boxsep=0.5mm,
      raster multicolumn=2, valign upper=center]
      \begin{equation*}
        \Psi_i(\xi_q) \mapsto \delta_{iq}
      \end{equation*}
    \end{tcolorbox}
    \begin{tcolorbox}[title=Enriched,left=0mm,right=0mm,top=0mm,bottom=0mm,boxsep=0.5mm,
      raster multicolumn=4, valign upper=center]
      \begin{equation*}
        {(\Phi \oplus \Psi)}_i(\xi_q) \mapsto
        \begin{cases}
          \Phi_i(\xi_q) &\text{$i < N_\Phi$}\\
          \Psi_{i-N_\Phi}(\xi_q) &\text{otherwise}
        \end{cases}
      \end{equation*}
    \end{tcolorbox}
    \begin{tcolorbox}[title=Tensor product,left=0mm,right=0mm,top=0mm,bottom=0mm,boxsep=0.5mm,
      raster multicolumn=3, valign upper=center]
      \begin{equation*}
        \Psi_{(i_1, \dots, i_n)}(\xi_{(q_1, \dots, q_n)}) \mapsto \prod_{j=1}^{n} \psi_{i_j}(\xi_{q_j})
      \end{equation*}
    \end{tcolorbox}
    \begin{tcolorbox}[title=$\mathbf{H}(\text{div}/\text{curl})$,left=0mm,right=0mm,top=0mm,bottom=0mm,boxsep=0.5mm,
      raster multicolumn=3, valign upper=center]
      \begin{equation*}
        \text{\texttt{HDiv}}(\Psi_i(\xi_q)) \mapsto
        \begin{cases}
          [-\Psi_i(\xi_q), 0]\\
          [0, \Psi_i(\xi_q)]
        \end{cases}
      \end{equation*}
    \end{tcolorbox}
  \end{tcbraster}
  \end{posterbox}

  \begin{posterbox}[name=rules, column=0, below=structure, span=3]{Compiler algorithms}
    Structure exploiting code generation implemented in TSFC, to
    obtain optimal complexity finite element kernels.

    \vspace{-\baselineskip}
    \subsection*{Delta cancellation}
    \vspace{-0.5em}
    \begin{equation*}
      \sum_{i_1, i_2} c_{i_1 i_2} \Psi_{i_1 q}
      \delta_{i_2 a} \mapsto \sum_{i_1} c_{i_1 a}\Psi_{i_1 q}
    \end{equation*}
    \vspace{-1.2\baselineskip}
    \subsubsection*{\dots across assignments}
    \vspace{-0.3em}
    \begin{equation*}
      A_{i_1 i_2 j_1 j_2} = \delta_{i_2 j_2} \sum_q w_q S_q \Psi_{i_1 q}\Psi_{j_1 q} \mapsto A_{i_1 i_2 j_1 i_2} = \sum_q w_q S_q \Psi_{i_1 q}\Psi_{j_1 q}
    \end{equation*}
    \vspace{-1.25\baselineskip}
    \subsection*{Sum factorisation}
    \vspace{-0.5em}
    An operation minimal contraction order is chosen by exhaustive search.
    \begin{align*}
      \sum_{q_1 q_2} D\Psi^{(1)}_{i_1 q_1}\Psi^{(2)}_{i_2 q_2}
      D\Psi^{(1)}_{j_1 q_1}\Psi^{(2)}_{j_2 q_2} P^{(1,1)}_{q_1 q_2}
      \mapsto \\
      \sum_{q_1} D\Psi^{(1)}_{i_1 q_1} D\Psi^{(1)}_{j_1 q_1} \sum_{q_2}  \Psi^{(2)}_{i_2 q_2}\Psi^{(2)}_{j_2 q_2} P^{(1,1)}_{q_1 q_2}
    \end{align*}
  \end{posterbox}

  \begin{posterbox}[name=flop-scaling, column=3, below=introduction, span=3]{Algorithmic complexity}
    \begin{center}
      FLOP/hexahedral cell computing
      \begin{equation*}
      \int\!\nabla u \cdot \nabla v\,\text{d}x.
    \end{equation*}

      \begin{tikzpicture}[scale=0.823]
        \begin{loglogaxis}[
          legend cell align=left,
          legend pos=outer north east,
          xlabel={Polynomial degree ($ \mathrm{n} $)},
          xtick={1,2,3,4,5,6,7,8,10,12,14,16,20,24,28,32},
          xlabel near ticks,
          xticklabels={1,2,3,4,,,,8,,,,16,,,,32},
          x tick label style={yshift=-0.15em},
          ylabel={FLOP/cell},
          ylabel near ticks,
          y tick label style={xshift=-0.15em},
          ]
          \addplot[densely dotted,mark=o,mark options={solid}, line width=1.5pt]
          table [x=degree, y=coffee, col sep=comma] {data/bipoisson_flops.csv};
          \addlegendentry{\texttt{coffee}}
          
          \addplot[densely dashed,domain=1:32] {36*(x+1)^9};
          \addlegendentry{$ \mathrm{(n+1)^9} $}
          
          \addplot[densely dotted,mark=*,mark options={solid}, line width=1.5pt]
          table [x=degree, y=spectral, col sep=comma] {data/bipoisson_flops.csv};
          \addlegendentry{\texttt{spectral}}
          
          \addplot[densely dashed,domain=1:32] {(x+1)^7};
          \addlegendentry{$ \mathrm{(n+1)^7} $}
          
          \addplot[densely dotted,mark=diamond*,mark options={solid}, line width=1.5pt]
          table [x=degree, y=gll_spectral, col sep=comma] {data/bipoisson_flops.csv};
          \addlegendentry{underintegration}
          
          \addplot[densely dashed,domain=1:32] {4*(x+1)^5};
          \addlegendentry{$ \mathrm{(n+1)^5} $}
        \end{loglogaxis}
      \end{tikzpicture}
    \end{center}
  \end{posterbox}

  \begin{posterbox}[name=curl-curl, column=3, below=flop-scaling,
    span=3]{DoF throughput}
    \begin{center}
      Throughput on a hexahedral mesh computing curl-curl action:
      \begin{equation*}
        \int_\Omega\!\nabla \times u \cdot \nabla \times v\,\text{d}x \quad u,v \in V \subset H(\text{curl}) 
      \end{equation*}
      
      \begin{tikzpicture}[scale=0.9]
      \begin{loglogaxis}[
        legend cell align=left,
        legend pos=outer north east,
        xlabel={Polynomial degree ($ \mathrm{n} $)},
        xtick={1,2,3,4,5,6,7,8,10,12,14,16,20,24,28,32},
        xlabel near ticks,
        xticklabels={1,2,3,4,,,,8,,,,16,,,,32},
        x tick label style={yshift=-0.15em},
        ylabel={DoF/$s$},
        ylabel near ticks,
        ]
        \addplot[densely dotted,mark=star,mark options={solid}, line width=1.5pt] table
        [x=degree, y=rate, col sep=comma] {data/curlcurl_spmv.csv};
        \addlegendentry{\textsc{MatVec}}

        \addplot[dotted,mark=*,mark options={solid,fill=gray}, line width=1.5pt] table
        [x=degree, y=rate, col sep=comma] {data/curlcurl_base.csv};
        \addlegendentry{FIAT/\texttt{coffee}}
        \pgfplotstableread{data/curlcurl.dat}{\data}
        \fitCurve{\data}{0}{\C}{\C * x^3 / (x+1)^6}
        \addlegendentry{$ \mathrm{n^3 / (n+1)^6} $}

        \addplot[dotted,mark=*,mark options={solid,fill}, line width=1.5pt] table
        [x=degree, y=rate, col sep=comma] {data/curlcurl_spectral.csv};
        \addlegendentry{\texttt{spectral}}
        \fitCurve{\data}{1}{\C}{\C * x^3 / (x+1)^4}
        \addlegendentry{$ \mathrm{n^3 / (n+1)^4} $}
      \end{loglogaxis}
    \end{tikzpicture}
  \end{center}
  \end{posterbox}

  \begin{posterbox}[name=future, column=0, below=rules,
    span=3]{Future work}
    \subsection*{Bernstein on simplices}
    \vspace{-0.5em}
      Sum factorisation on simplex elements using Bernstein basis.

      Requires mostly similar technology, but now tensor products are
      ragged.
      \vspace{-\baselineskip}
      \subsection*{Local tensor algebra}
      \vspace{-0.5em}
      Fast tensor algebra for hybridisation methods.  Requires
      assembling (or providing action of) element tensors with
      structure $D - C A^{-1} B$.

      Should also exploit tensor product identities where possible,
      e.g.~$\big[\otimes_i A_i\big]^{-1} = \otimes_i A_i^{-1}$.
      \vspace{-\baselineskip}
      \subsection*{Vectorisation}
      \vspace{-0.5em}
      Working in progress cross-element vectorisation schemes.  Have
      low-complexity algorithms, need to do the performance analysis
      to ensure high performance too.
  \end{posterbox}    
  \begin{posterbox}[name=try, column=3, below=curl-curl,
    span=3]{Try it}
    \subsection*{More details}
    \vspace{-0.5em}
    Homolya, Kirby, Ham, \emph{Exposing and exploiting structure:
      optimal code generation for high-order finite element
      methods}. To appear in ACM TOMS (2018).
    \texttt{arXiv:\,1711.02473\,[cs.MS]}.
    \vspace{-\baselineskip}
    \subsection*{Code}
    \vspace{-0.5em}
    FInAT: \url{https://github.com/finat/finat/}\\
    TSFC: \url{https://github.com/firedrakeproject/tsfc/}\\
    Firedrake: \url{https://www.firedrakeproject.org/}

    \vspace{\baselineskip}
    \begin{minipage}{0.167\textwidth}
      \begin{flushleft}
        \includegraphics[align=c,height=2.5em]{epsrc-logo}
      \end{flushleft}
    \end{minipage}
    \hfill
    \begin{minipage}{0.167\textwidth}
      \begin{center}
        \includegraphics[align=c,height=2.5em]{nerc-logo}
      \end{center}
    \end{minipage}
    \hfill
    \begin{minipage}{0.167\textwidth}
      \begin{flushright}
        \includegraphics[align=c,height=3.5em]{nsf-logo}
      \end{flushright}
    \end{minipage}
  \end{posterbox}
  \begin{posterbox}[name=acknowledgements,
    column=0,span=6,below=future,boxheaderheight=2em]{Support}
    The Grantham Institute; the National Science Foundation [award number 1525697]; the Natural Environment
    Research Council [grant number NE/K008951/1]; the Engineering and
    Physical Sciences Research Council [grant number EP/M011054/1].
  \end{posterbox}
\end{poster}
\end{document}
